\chapter{Ergodic theorems}

Recall the classical law of large numbers: if \(X_0,X_1,\dots\) are i.i.d.\ integrable random variables, then
\[
\frac1n\sum_{k=0}^{n-1}X_k \xrightarrow[n\to\infty]{} \E[X_0]
\quad \text{a.s.}
\]
The key assumption here is \emph{independence}. In many situations, however, the sequence \((X_k)\) is not independent:
for instance time series, dynamical systems, or Markov chains. Still, one often expects that \emph{time averages}
\[
\frac1n\sum_{k=0}^{n-1}X_k(\omega)
\]
converge to a deterministic quantity that describes the \emph{typical long-time behavior} of the system.
Ergodic theory provides conditions under which such a replacement is valid: we may ``skip independence'' but keep
a suitable invariance (stationarity) and an irreducibility-type property (ergodicity).

\medskip

\noindent
\textbf{Guiding principle.}
Ergodic theorems identify situations in which
\[
\text{time average} \ \approx\ \text{space average}.
\]
The time average is computed along a single trajectory \(\omega \mapsto (X_k(\omega))\),
whereas the space average is an expectation with respect to the underlying probability measure \(\Pp\).

%========================================================
\section{Setting and notation}
Let
\[
E=\{1,\dots,\ell\},
\qquad (X_t)_{t\in\N} \ \text{a Markov chain on } E.
\]
We write the law (distribution) of \(X_t\) as the row vector
\[
\alpha_t = (\alpha_t^1,\dots,\alpha_t^\ell),
\qquad \alpha_t^i := \Pp(X_t=i).
\]
Let \(P=(p_{ij})_{i,j\in E}\) be the transition matrix, i.e.
\[
p_{ij} := \Pp(X_{t+1}=j \mid X_t=i).
\]
Then the distribution evolves as
\[
\alpha_{t+1} = \alpha_t P,
\qquad\text{hence}\qquad
\alpha_t = \alpha_0 P^t.
\]

\paragraph{Stationary distributions.}
In ergodic theory one is interested in \emph{long--time averages} along trajectories and their relation to
\emph{space averages} with respect to an invariant law. For a Markov chain with transition matrix $P$,
the distribution at time $t$ evolves by
\[
\alpha_{t+1}=\alpha_t P.
\]
Hence we need a notion of a probability law that is \emph{preserved by the dynamics} (i.e.\ invariant under one step of the chain):
a probability vector $\pi$ such that
\[
\pi=\pi P.
\]
If the chain starts in such a law, $\alpha_0=\pi$, then $\alpha_t=\pi$ for every $t\ge 0$, so the process is
\emph{stationary in time} and time shifts do not change its distribution. This is the natural probabilistic analogue
of an \emph{invariant measure} in dynamical systems.

Stationary distributions are crucial because they provide the canonical candidate for the ``equilibrium'' law
governing asymptotic behavior: under suitable irreducibility/aperiodicity assumptions, one expects
\[
\alpha_t \;\to \pi \quad (t\to\infty),
\]
and, more importantly for ergodic theorems, they are the measures with respect to which time averages
\[
\frac1n\sum_{t=0}^{n-1} f(X_t)
\]
can be compared to the corresponding space average $\sum_{i=1}^\ell f(i)\,\pi^i$.
Thus the stationary distributions singles out exactly those laws for which the dynamics does not ``drift'' in distribution,
making them the correct reference measures for formulating and proving ergodic results.


\begin{definition}[Stationary distributions]

A probability vector \(\pi = (\pi^1,\dots,\pi^\ell)\) is called a \emph{stationary distribution} if
\[
\pi = \pi P
\qquad\textrm{or equivalently}\qquad
\pi^T = P^T \pi^T.
\]
In particular, if \(\alpha_0=\pi\), then \(\alpha_t=\pi\) for all \(t\), i.e.\ the chain is stationary in time.

\end{definition}

\paragraph{Ergodic Markov chains.}
A central goal of ergodic theorems is to formalize the idea that a system \emph{settles into equilibrium} and
that long--time observations become independent of the initial condition. For a Markov chain, the most direct
way to express ``equilibrium in the long run'' is in terms of the $t$--step transition probabilities
\[
p_{ij}^{(t)}=\Pp(X_t=j\mid X_0=i)=(P^t)_{ij}.
\]
If, as $t\to\infty$, these probabilities converge to a limit $\pi_j$ that is \emph{independent of the starting state $i$},
then the chain has \emph{forgotten its initial state}: regardless of where we start, the law of $X_t$ converges to the
same distribution $\pi$. This is precisely the type of asymptotic stability one needs in order to speak about a
\emph{unique equilibrium law} and to justify replacing time averages along a single trajectory by expectations under $\pi$.

\begin{definition}[Ergodic Markov Chains]

For \(t\ge 0\), denote the \(t\)-step transition probabilities by
\[
p_{ij}^{(t)} := (P^t)_{ij} = \Pp(X_t=j \mid X_0=i).
\]
The chain is called \emph{ergodic} (in the sense used in the notes) if the limit
\[
\pi_j := \lim_{t\to\infty} p_{ij}^{(t)}
\]
exists, is \emph{independent of \(i\)}, and satisfies
\[
\pi_j>0,
\qquad
\sum_{j\in E}\pi_j = 1.
\]
In this case, \(\pi\) is necessarily the stationary distribution and the chain forgets its initial state in the long run.
    
\end{definition}

In short, ergodicity packages the two key asymptotic features---\emph{convergence to equilibrium} and \emph{loss of memory of the initial state}---into a single clean condition.

\begin{remark}
    
The additional requirements $\pi_j>0$ and $\sum_{j}\pi_j=1$ ensure that the limit is a genuine probability distribution
with full support (no state is ``asymptotically ignored''), matching the intuitive picture that in equilibrium the chain
keeps visiting every state. Moreover, the definition is tailored so that, once the limit exists,
$\pi$ must be stationary: since each row of $P^t$ converges to $\pi$, one obtains $\pi=\pi P$, and hence $\pi$ is the
invariant (equilibrium) distribution. 

\end{remark}

%========================================================
\section{Stationary stochastic processes}

For Markov chains we singled out a \emph{stationary distribution} $\pi$ as an equilibrium law satisfying
$\pi=\pi P$, so that if $X_0\sim\pi$ then the one--time marginal distribution does not change with $t$.
We then called a chain \emph{ergodic} if, regardless of the initial state, the $t$--step transition
probabilities converge to a limit $\pi$ independent of the start, i.e.\ the chain \emph{forgets its past} and
approaches equilibrium. Both notions are, however, tied to a particular model class (Markov chains) and are phrased
mainly in terms of \emph{marginals} such as $\mathcal{L}(X_t)$. 

In ergodic theory one wants a concept that applies to
\emph{arbitrary} stochastic processes and captures invariance not only of single-time laws but of the \emph{entire}
finite-dimensional distributions. This leads to the definition of a stationary process: the joint law of
$(X_t)_{t\in I}$ is invariant under time shifts, meaning that for every shift $s$ the shifted process
$(X_{t+s})_{t\in I}$ has the same distribution as $(X_t)_{t\in I}$. In particular, all marginals are time-homogeneous,
and the process has no distinguished ``origin of time''---exactly the setting in which one can hope to relate
time averages along trajectories to expectations under the common (invariant) law.



\begin{definition}

Let \(I\) be an index set (typically \(I=\N\) or \(I=\Z\)). A stochastic process
\[
X=(X_t)_{t\in I} \quad \text{on } (\Omega,\A,\Pp)
\]
is called \emph{stationary} if it is invariant under time shifts:
\[
(X_{t+s})_{t\in I} \ \stackrel{d}{=} \ (X_t)_{t\in I}
\qquad \forall s\in I.
\]

Equivalently, for any \(t_1,\dots,t_m\in I\),
\[
(X_{t_1},\dots,X_{t_m}) \ \stackrel{d}{=} \ (X_{t_1+s},\dots,X_{t_m+s})
\qquad \forall s\in I.
\]
    
\end{definition}

\begin{remark}
To use the shift operator as an \emph{invertible} time-translation (a standard setup in ergodic theory), we can extend a stationary one-sided process on $\N$ to a two-sided stationary process on $\Z$.

\begin{lemma}[Lemma (two-sided extension)]
Every stationary process \(X=(X_t)_{t\in\N}\) can be extended to a stationary process
\[
\bar X=(\bar X_t)_{t\in\Z}.
\]
(Conceptually: stationarity allows one to consistently define the joint laws backwards in time.)

\end{lemma}
    
\end{remark}



%========================================================
\section{Measure-preserving dynamical systems}

Up to now, ``stationarity'' meant that the probabilistic description of a process does not change under a time shift.
The next step in ergodic theory is to turn this symmetry into a \emph{dynamical system} by introducing a measurable
map $T:\Omega\to\Omega$ that represents ``one step forward in time'' (for a process, $T$ is typically the left shift on
the path space). Once we have $T$, we can ask a basic structural question:

\medskip
\emph{Which properties of the system can be detected from infinitely long observation along the orbit
$\omega, T\omega, T^2\omega,\dots$?}

\medskip
The answer is encoded by the \emph{invariant $\sigma$-algebra} $\mathcal{I}$: an event $A$ is $T$-invariant if it is
unaffected by moving one step forward in time, so membership in $A$ depends only on the ``tail'' behavior of the orbit.
These events represent the \emph{persistent information} that the dynamics cannot wash out. Measure-preservation is the
minimal compatibility condition ensuring that shifting time does not change probabilities, i.e.\ $T$ does not distort
the underlying law.

Finally, \emph{ergodicity} expresses the idea that the system has \emph{no nontrivial persistent information}:
the only invariant events have probability $0$ or $1$. Equivalently, the only random variables that remain unchanged
under the time shift (i.e.\ $X=X\circ T$ a.s.) are constants. This viewpoint is the bridge to ergodic theorems:
it formalizes ``the system cannot decompose into different invariant regimes'', and it prepares the key conclusion that
long-time averages cannot depend on hidden invariant features of the initial state.

\paragraph{Invariant events and the invariant \(\sigma\)-algebra.}
Let \((\Omega,\A,\Pp)\) be a probability space and \(T:\Omega\to\Omega\) be measurable. The \emph{invariant $\sigma$-algebra} of $T$ is represents the collection of all events that are unchanged under the dynamics.

\begin{definition}[Invariant and quasi-invariant sets]
A set \(A\in\A\) is called
\begin{itemize}
\item \emph{\(T\)-invariant} if \(T^{-1}(A)=A\),
\item \emph{quasi-invariant} if \(\mathbf{1}_{T^{-1}(A)}=\mathbf{1}_A\) \(\Pp\)-a.s.
\end{itemize}
The collection of invariant sets forms a \(\sigma\)-algebra:
\[
\I := \{A\in\A \,:\, T^{-1}(A)=A\},
\]
called the \emph{invariant \(\sigma\)-algebra}.
\end{definition}

\paragraph{Measure-preservation.} We are interesting in systems $T$ that advances in time but does not distort probabilities, so the dynamics is compatible with the reference law $\Pp$. 
This is exactly the condition that guarantees that time shifts leave distributions (and hence expectations) unchanged, making long-time averages comparable to $\Pp$-averages. 
It is therefore the natural framework in which ergodic statements (``time averages = space averages'') can even be formulated.

\begin{definition}[Measure-preserving system]
The map \(T\) is called \emph{measure-preserving} if
\[
\Pp(T^{-1}(A)) = \Pp(A)
\qquad \forall A\in\A.
\]
In this case, \((\Omega,\A,\Pp,T)\) is called a \emph{measure-preserving dynamical system}.
\end{definition}

\paragraph{Ergodicity.}
The invariant $\sigma$-algebra $\mathcal I$ collects exactly those events whose occurrence is unchanged by the time-shift,
so they represent the ``information'' about the system that persists forever along an orbit.
Requiring $\mathcal I$ to be $\Pp$-trivial means there are no nontrivial invariant regimes of positive probability:
every invariant event either happens almost surely or almost never.
Equivalently, the only random variables that are unchanged by the dynamics ($X=X\circ T$ a.s.) are constants, which captures
the idea that long-time behavior cannot depend on hidden time-invariant features of the initial state.


\begin{definition}[Ergodicity]
A measure-preserving system \((\Omega,\A,\Pp,T)\) is called \emph{ergodic} if the invariant \(\sigma\)-algebra \(\I\)
is \(\Pp\)-trivial, i.e.
\[
A\in\I \ \Longrightarrow\ \Pp(A)\in\{0,1\}.
\]
\end{definition}

If the system is ergodic, the invariant information is trivial, hence there exist observable which are almost surely constant. This is the key mechanism behind ``time averages cannot depend on the initial condition''. The next lemma identifies $\mathcal I$-measurable random variables as exactly those observables that are unchanged by the dynamics, i.e.\ $X(\omega)=X(T\omega)$ almost surely, so they represent the ``time-invariant information'' of the system. 


\begin{lemma}[Invariant random variables]\label{lem:invariant_rv}
Let \((\Omega,\A,\Pp,T)\) be measure-preserving and \(X:\Omega\to\R\) a random variable.
\begin{enumerate}[label=(\alph*)]
\item \(X\) is \(\I\)-measurable \(\iff\) \(X = X\circ T\) \(\Pp\)-a.s.
\item If the system is ergodic, then any \(\I\)-measurable \(X\) is \(\Pp\)-a.s.\ constant.
\end{enumerate}
\end{lemma}

\begin{proof}
(a) ``\(\Rightarrow\)'' follows since \(\I\)-measurability means invariance of all level sets, which is equivalent to
\(X(\omega)=X(T\omega)\) a.s.  ``\(\Leftarrow\)'' is immediate from the definition of \(\I\).
(b) If \(X\) is \(\I\)-measurable, then \(\{X\le c\}\in\I\) for every \(c\in\R\). Ergodicity forces
\(\Pp(X\le c)\in\{0,1\}\), which implies that \(X\) is a.s.\ constant.
\end{proof}

%========================================================
\section{The canonical shift system of a process}

Stationarity is a statement about \emph{time shifts},i.e., looking at a process from time $0$ or from time $s$ should make no
difference in distribution. To exploit this symmetry systematically, it is convenient to encode \emph{one whole
realization} of the process as a single point $\omega$ in a new space, the \emph{path space} $E^\N$. On this space,
``moving one step forward in time'' becomes an honest deterministic map $T$ (the left shift), and iterating $T$
corresponds exactly to observing later and later parts of the same trajectory.

\paragraph{Path space and the shift map.}
Let \((E,\E)\) be a measurable space and consider the path space
\[
\bigl(E^\N,\,\E^{\otimes \N}\bigr).
\]
Define the (left) shift
\[
T:E^\N \to E^\N,
\qquad
T((x_0,x_1,x_2,\dots)) := (x_1,x_2,x_3,\dots).
\]
Then
\[
T^n((x_0,x_1,x_2,\dots)) = (x_n,x_{n+1},x_{n+2},\dots).
\]

\medskip
Once we view trajectories as points, the random variables of the process become the simplest possible observables:
the coordinate maps $X_n(\omega)=\omega_n$. The identity $X_n=X_0\circ T^n$ then says in one line that \emph{the entire
process is generated by repeatedly applying the same time-shift to the same observable} $X_0$.
This is the key bridge between probability and dynamics; questions about long-time behavior of $(X_n)$ turn into
questions about iterates $T^n$.

\paragraph{Coordinate process.}
Define the coordinate maps
\[
X_n:E^\N \to E,
\qquad
X_n((x_k)_{k\in\N}) := x_n.
\]
They satisfy the fundamental relation
\[
X_n = X_0\circ T^n
\qquad (n\in\N).
\]

\medskip

Finally, the stationarity assumption becomes a structural property of the shift system. The law $\Pp$ on path space is invariant under $T$, i.e.\ $\Pp(T^{-1}A)=\Pp(A)$. In other words, a stationary process
is precisely a \emph{measure-preserving dynamical system} on its canonical path space.
This reformulation is useful because it lets us apply the general machinery of ergodic theory (invariant
$\sigma$-algebra, ergodicity, ergodic theorems) directly to stochastic processes by studying the single map $T$.

\paragraph{Stationarity as measure-preservation}
Let \(\Pp\) be the law of a process \((X_n)_{n\in\N}\) on \(E^\N\). Then the process is stationary if and only if
\[
\Pp(T^{-1}(A))=\Pp(A)\qquad\forall A\in \E^{\otimes \N},
\]
i.e.\ \((E^\N,\E^{\otimes \N},\Pp,T)\) is a measure-preserving dynamical system.
Thus, stationary processes can be studied through measure-preserving dynamics.

%========================================================
\section{Ergodic theorem setting and time averages}

Having recast stationarity as measure-preservation, we can now formulate the central question of ergodic theory in its
cleanest form: \emph{what does one learn about the law $\Pp$ by observing a single trajectory for a long time?}
A measurable function $f:\Omega\to\R$ plays the role of an \emph{observable} (a quantity we can measure from the state
$\omega$), and iterating the dynamics produces the sequence of observations
\[
X_n(\omega)=f(T^n\omega).
\]
Thus the process is generated by repeatedly applying the same time shift to the same observable, and measure-preservation
guarantees that these observations form a stationary process.

\medskip
The objects of interest are the \emph{time averages}
\[
S_n(\omega)=\frac1n\sum_{k=0}^{n-1} f(T^k\omega),
\]
which represent empirical averages along one orbit.
The purpose of this section is to understand the limit of $S_n(\omega)$ as $n\to\infty$ and, in particular, to connect it
to the corresponding \emph{space average} $\int_\Omega f\,d\Pp$; this is the precise mathematical meaning of the slogan
``time averages equal ensemble averages'' under suitable ergodicity assumptions.


Let \((\Omega,\A,\Pp,T)\) be a measure-preserving dynamical system and let
\[
X_0=f:\Omega\to\R
\quad\text{be measurable}.
\]
Define
\[
X_n(\omega) := f(T^n\omega) = X_0(T^n\omega),
\qquad n\in\N.
\]
Then \((X_n)_{n\in\N}\) is a real-valued stationary process.

The associated time averages are
\[
S_n(\omega)
:= \frac1n\sum_{k=0}^{n-1}X_k(\omega)
= \frac1n\sum_{k=0}^{n-1} f(T^k\omega),
\qquad n\ge 1,
\]
with \(T^0=\textrm{id}\).

%========================================================
For proving the main result of this chapter, we will need the lemma below. 

\begin{lemma}[Hopf maximal ergodic lemma]
Assume \(X_0\in L^1(\Pp)\). Define the running maximum
\[
M_n := \max\{0,S_1,S_2,\dots,S_n\}.
\]
Then
\[
\E\bigl[ X_0\, \mathbf{1}_{\{M_n>0\}} \bigr] \ge 0
\qquad \forall n\in\N.
\]
\end{lemma}

\begin{proof}
Let $f:=X_0\in L^1(\Pp)$ and define the (non-averaged) partial sums
\[
\Sigma_0:=0,
\qquad
\Sigma_k(\omega):=\sum_{j=0}^{k-1} f(T^j\omega),\quad k\ge 1.
\]
Then $S_k(\omega)=\frac{1}{k}\Sigma_k(\omega)$, so
\[
M_n(\omega)>0
\ \Longleftrightarrow\
\exists\,k\in\{1,\dots,n\}\ \text{with}\ S_k(\omega)>0
\ \Longleftrightarrow\
\exists\,k\in\{1,\dots,n\}\ \text{with}\ \Sigma_k(\omega)>0.
\]
Hence, if we set
\[
m_n(\omega):=\max\{\Sigma_0(\omega),\Sigma_1(\omega),\dots,\Sigma_n(\omega)\},
\qquad
A:=\{M_n>0\}=\{m_n>0\},
\]
it suffices to prove $\E[f\,\mathbf 1_A]\ge 0$.

\medskip
\Step{1: Choose a first time of reaching the (positive) maximum}
On $A$ define
\[
\tau(\omega):=\min\{1\le k\le n:\ \Sigma_k(\omega)=m_n(\omega)\}.
\]
Then $\tau(\omega)\in\{1,\dots,n\}$ and $\Sigma_{\tau(\omega)}(\omega)=m_n(\omega)>0$.

\medskip
\Step{2: Compare $m_n(\omega)$ and $m_n(T\omega)$}
For $\omega\in A$, we have $\tau(\omega)-1\in\{0,\dots,n-1\}$, hence by definition of $m_n(T\omega)$,
\[
m_n(T\omega)\ \ge\ \Sigma_{\tau(\omega)-1}(T\omega).
\]
But
\[
\Sigma_{\tau-1}(T\omega)
=\sum_{j=0}^{\tau-2} f(T^{j+1}\omega)
=\sum_{j=1}^{\tau-1} f(T^{j}\omega)
=\Sigma_\tau(\omega)-f(\omega)
=m_n(\omega)-f(\omega).
\]
Therefore, for $\omega\in A$,
\[
m_n(T\omega)\ \ge\ m_n(\omega)-f(\omega)
\quad\Longrightarrow\quad
f(\omega)\ \ge\ m_n(\omega)-m_n(T\omega).
\]

\medskip
\Step{3: Extend the inequality to all $\omega$}
If $\omega\notin A$, then $m_n(\omega)=0$ and $\mathbf 1_A(\omega)=0$, so
\[
f(\omega)\mathbf 1_A(\omega)=0\ \ge\ m_n(\omega)-m_n(T\omega)
\]
since the right-hand side is $\le 0$. Hence, for all $\omega\in\Omega$,
\[
f(\omega)\mathbf 1_A(\omega)\ \ge\ m_n(\omega)-m_n(T\omega).
\]

\medskip
\Step{4: Integrate and use measure-preservation}
Integrating and using that $T$ is measure-preserving (so $\E[m_n\circ T]=\E[m_n]$), we obtain
\[
\E\bigl[f\,\mathbf 1_A\bigr]
\ \ge\
\E[m_n]-\E[m_n\circ T]
\ =\ 0.
\]
Since $A=\{M_n>0\}$, this is exactly
\[
\E\bigl[X_0\,\mathbf 1_{\{M_n>0\}}\bigr]\ge 0,
\]
as claimed.
\end{proof}


\noindent
\emph{Remark}
The event \(\{M_n>0\}\) means that at some time \(1\le k\le n\) the average \(S_k\) becomes positive.
The lemma states that the contribution of \(X_0\) on this event must have nonnegative expectation.
This inequality is the key estimate that leads to almost sure convergence of \(S_n\).

%========================================================
\section{Birkhoff's individual ergodic theorem}

We are now ready to state the main result that justifies \emph{``time averages equal space averages''}
for systems generated by a measure-preserving time shift. In the setting of a dynamical system
$(\Omega,\mathcal A,\Pp,T)$ and an observable $f\in L^1(\Pp)$, we observe the dependent sequence
$X_n=f\circ T^n$ along a single trajectory and form the empirical (time) averages
\[
S_n(\omega)=\frac1n\sum_{k=0}^{n-1} f(T^k\omega).
\]
The question is whether these averages converge, and if so, what they converge to and how the answer depends on
possible invariant structure (encoded by the invariant $\sigma$-algebra $\mathcal I$).

\medskip
Birkhoff's individual ergodic theorem gives a complete answer: for every integrable observable the time averages
converge almost surely (and in $L^1$), and the limit is the $\mathcal I$-measurable part of $f$, namely
the conditional expectation $\E[f\mid\mathcal I]$. In particular, when the system is ergodic (so $\mathcal I$ is trivial),
the limit becomes the constant $\E[f]$, showing that a single typical long trajectory is representative of the whole
distribution.


\begin{theorem}[Birkhoff (individual ergodic theorem)]
Let \((\Omega,\A,\Pp,T)\) be measure-preserving and \(f\in L^1(\Pp)\). Define \(X_n=f\circ T^n\) and
\[
S_n=\frac1n\sum_{k=0}^{n-1}f\circ T^k.
\]
Then \(S_n\) converges \(\Pp\)-a.s.\ and in \(L^1(\Pp)\) to the conditional expectation of \(f\) w.r.t.\ the invariant \(\sigma\)-algebra \(\I\):
\[
\lim_{n\to\infty} S_n
= \E[f \mid \I]
= \E[X_0 \mid \I]
\qquad \Pp\text{-a.s.}
\]
If the system is ergodic, then \(\I\) is trivial and hence
\[
\lim_{n\to\infty} S_n = \E[f] = \E[X_0]
\qquad \Pp\text{-a.s.}
\]
\end{theorem}

\begin{proof}
Write $f:=X_0\in L^1(\Pp)$ and $X_n=f\circ T^n$.  Set
\[
S_n(\omega)=\frac1n\sum_{k=0}^{n-1} f(T^k\omega),\qquad
\overline f(\omega):=\limsup_{n\to\infty} S_n(\omega),\qquad
\underline f(\omega):=\liminf_{n\to\infty} S_n(\omega).
\]

\medskip
\Step{1: $\overline f$ and $\underline f$ are $T$-invariant}
Using
\[
S_n(T\omega)
=\frac1n\sum_{k=0}^{n-1} f(T^{k+1}\omega)
=\frac1n\sum_{k=1}^{n} f(T^{k}\omega)
=S_n(\omega)+\frac{f(T^n\omega)-f(\omega)}{n},
\]
we get $S_n(T\omega)-S_n(\omega)\to 0$ for a.e.\ $\omega$ (since $f(T^n\omega)$ is finite a.e.).
Hence
\[
\overline f\circ T=\overline f,\qquad \underline f\circ T=\underline f \qquad \Pp\text{-a.s.}
\]
In particular, $\overline f$ and $\underline f$ are $\mathcal I$-measurable (by Lemma~\ref{lem:invariant_rv}(a)).

\medskip
\Step{2: A key inequality from Hopf's maximal lemma}
Fix any $\mathcal I$-measurable (equivalently $T$-invariant) function $g\in L^1(\Pp)$ and apply the
Hopf maximal ergodic lemma to $h:=f-g\in L^1(\Pp)$.
Denote the corresponding time averages by
\[
S_n(h)=\frac1n\sum_{k=0}^{n-1} (f-g)\circ T^k.
\]
Since $g\circ T^k=g$ a.s., we have the identity
\[
S_n(h)=S_n(f)-g\qquad\text{a.s. for every }n.
\]
Define the running maximum $M_n(h):=\max\{0,S_1(h),\dots,S_n(h)\}$ and the event
\[
A_n:=\{M_n(h)>0\}=\Bigl\{\max_{1\le k\le n}\bigl(S_k(f)-g\bigr)>0\Bigr\}.
\]
Hopf's lemma yields, for every $n$,
\[
\E\bigl[(f-g)\,\mathbf 1_{A_n}\bigr]\ge 0.
\]
Since $A_n\uparrow A:=\bigcup_{n\ge 1}A_n=\{\sup_{k\ge 1}(S_k(f)-g)>0\}$, by monotone convergence,
\[
\E\bigl[(f-g)\,\mathbf 1_{A}\bigr]\ge 0.
\]
Now observe that $\{\overline f>g\}\subset A$: if $\overline f(\omega)>g(\omega)$, then for some $k$ we must have
$S_k(f)(\omega)>g(\omega)$, hence $\sup_k(S_k(f)-g)(\omega)>0$.
Therefore, with $B:=\{\overline f>g\}$,
\begin{equation}\label{eq:keyHopfIneq}
\E\bigl[(f-g)\,\mathbf 1_{B}\bigr]\ge 0
\qquad\Longleftrightarrow\qquad
\int_{B} f\,d\Pp \ \ge\ \int_{B} g\,d\Pp.
\end{equation}

\medskip
\Step{3: Minimality of $\overline f$ among invariant majorants}
Let $g\in L^1(\Pp)$ be $\mathcal I$-measurable and satisfy $g\ge \overline f$ a.s.
Then $B=\{\overline f>g\}$ has $\Pp(B)=0$, and \eqref{eq:keyHopfIneq} gives
\[
\int_\Omega f\,d\Pp \ \le\ \int_\Omega g\,d\Pp.
\]
(Indeed, $\int_B f\ge \int_B g$ is trivial when $\Pp(B)=0$, and on $B^c$ we have $f\le g$ only in average; the
argument is standard: apply \eqref{eq:keyHopfIneq} to $g_\varepsilon:=g-\varepsilon$ and let $\varepsilon\downarrow 0$
to deduce $\int f\le \int g$.)

In particular, taking $g:=\overline f$ itself shows $\overline f\in L^1(\Pp)$ and the inequality is sharp.

\medskip
\Step{4: Maximality of $\underline f$ among invariant minorants}
Apply Step~2 to $-f$ and $-g$ (or repeat the argument with $\underline f$).
Equivalently: for any $\mathcal I$-measurable $g\in L^1(\Pp)$ with $g\le \underline f$ a.s., one obtains
\[
\int_\Omega f\,d\Pp \ \ge\ \int_\Omega g\,d\Pp.
\]
In particular, taking $g:=\underline f$ shows $\underline f\in L^1(\Pp)$.

\medskip
\Step{5: Conclude $\overline f=\underline f$ a.s.\ and define the limit}
From Steps~3--4 with $g=\overline f$ and $g=\underline f$ we obtain
\[
\int_\Omega \overline f\,d\Pp \ \le\ \int_\Omega f\,d\Pp \ \le\ \int_\Omega \underline f\,d\Pp.
\]
But always $\underline f\le \overline f$ a.s., hence
\[
0\ \le\ \int_\Omega (\overline f-\underline f)\,d\Pp
\ \le\ \int_\Omega \overline f\,d\Pp-\int_\Omega \underline f\,d\Pp
\ \le\ 0,
\]
so $\int (\overline f-\underline f)\,d\Pp=0$ and therefore $\overline f=\underline f$ a.s.
Thus $S_n(\omega)$ converges a.s.; denote the limit by
\[
f^*(\omega):=\lim_{n\to\infty} S_n(\omega).
\]
By Step~1, $f^*$ is $T$-invariant, hence $\mathcal I$-measurable.

\medskip
\Step{6: Identify the limit as $\E[f\mid\mathcal I]$}
Let $A\in\mathcal I$ be invariant. Consider the induced measure-preserving system on $A$ and apply the Hopf inequality
to $f$ restricted to $A$ (equivalently, apply the previous steps to the function $f\,\mathbf 1_A$ and use that
$\mathbf 1_A\circ T=\mathbf 1_A$ a.s.). One obtains
\[
\int_A f^*\,d\Pp=\int_A f\,d\Pp
\qquad\forall A\in\mathcal I.
\]
Since $f^*$ is $\mathcal I$-measurable, this is exactly the defining property of the conditional expectation:
\[
f^*=\E[f\mid\mathcal I]\qquad \Pp\text{-a.s.}
\]
(Indeed, $Y=\E[f\mid\mathcal I]$ is the unique $\mathcal I$-measurable random variable satisfying
$\int_A Y\,d\Pp=\int_A f\,d\Pp$ for all $A\in\mathcal I$.)

\medskip
\Step{ 7: $L^1$-convergence}
First assume $f$ is bounded: $|f|\le C$. Then $|S_n|\le C$ and $S_n\to f^*$ a.s.,
so dominated convergence yields $S_n\to f^*$ in $L^1(\Pp)$.

For general $f\in L^1(\Pp)$, define truncations $f_m:=(-m)\vee f\wedge m\in L^\infty$.
By the bounded case, $S_n(f_m)\to \E[f_m\mid\mathcal I]$ in $L^1$.
Moreover, the averaging operator is an $L^1$-contraction:
\[
\|S_n(f-f_m)\|_{L^1}\le \frac1n\sum_{k=0}^{n-1}\| (f-f_m)\circ T^k\|_{L^1}
=\|f-f_m\|_{L^1},
\]
using measure-preservation, and conditional expectation is also an $L^1$-contraction:
\[
\|\E[f-f_m\mid\mathcal I]\|_{L^1}\le \|f-f_m\|_{L^1}.
\]
Hence, for every $n,m$,
\[
\|S_n(f)-\E[f\mid\mathcal I]\|_{L^1}
\le \|S_n(f)-S_n(f_m)\|_{L^1}
+\|S_n(f_m)-\E[f_m\mid\mathcal I]\|_{L^1}
+\|\E[f_m-f\mid\mathcal I]\|_{L^1}
\]
\[
\le 2\|f-f_m\|_{L^1}+\|S_n(f_m)-\E[f_m\mid\mathcal I]\|_{L^1}.
\]
Let $n\to\infty$ (for fixed $m$) and then $m\to\infty$ (since $f_m\to f$ in $L^1$) to conclude
\[
S_n(f)\to \E[f\mid\mathcal I]\qquad\text{in }L^1(\Pp).
\]

\medskip
\Step{8: The ergodic case}
If the system is ergodic, $\mathcal I$ is $\Pp$-trivial, so $\E[f\mid\mathcal I]$ is a.s.\ constant and equals $\E[f]$.
Thus $\lim_{n\to\infty}S_n=\E[f]$ a.s.\ (and in $L^1$).
\end{proof}


\noindent
\textbf{Interpretation.}
The limit \(\E[f\mid\I]\) is the ``best invariant approximation'' of \(f\): it depends only on the part of the
probability space that cannot be changed by shifting along trajectories.
In an ergodic system, there are no nontrivial invariant events, so the limit must be a constant; thus
time averages coincide with the global expectation.

\medskip

\noindent
\textbf{Takeaway.}
Birkhoff's theorem is the law of large numbers for dependent data produced by a measure-preserving system:
stationarity replaces i.i.d., and ergodicity replaces the absence of hidden invariant structure.


